\section{Introduction}

A reliable initial velocity model is important for seismic imaging and waveform-based inversion because velocity controls travel-time kinematics and the positioning and focusing of subsurface structures \citep{stolt1978migration,tarantola1984inversion,virieux2009overview}. Full-waveform inversion can recover detailed velocity structure, but its optimization remains sensitive to the starting model, acquisition coverage, missing low-frequency information, and cycle skipping \citep{bunks1995multiscale,virieux2009overview}. Initial velocity model building therefore remains relevant even when sophisticated downstream inversion tools are available.

Data-driven velocity reconstruction provides a complementary route by amortizing the mapping from observations to a velocity model \citep{yang2019deeplearning,wu2020inversionnet,deng2022openfwi}. The task studied here is specifically \emph{velocity-model reconstruction when processed or interpreted constraints are already available}: RMS velocity, a post-stack migrated image, interpreted horizons, and sparse well velocities. These constraints differ in physical meaning, sampling domain, density, and information content.

A common learned strategy is to encode the available observations, fuse their features, and assign the complete velocity field to one downstream estimator. The estimator may be a deterministic convolutional network, an adversarial generator, a latent model, or a diffusion/flow-based model. This design can exploit multimodal information, but it leaves reconstruction responsibility implicit: the same estimator must recover both a large-scale velocity background and the remaining structural correction. Increasing the capacity of a fused full-field model does not by itself specify how those responsibilities should be organized.

We formulate the problem differently. Every available modality can contribute to both learned condition roles, but the downstream estimators have explicit responsibilities. A background-oriented interface drives deterministic background estimation, while a structure-oriented interface guides the remaining correction. The predicted background is treated as a first-class inference variable rather than an intermediate feature: it defines the residual target, provides context to the residual generator, and is added back to the generated correction.

This leads to Background-Guided Residual Flow Matching (\method). Given the multimodal observations, the method first predicts a background field $\hat B$. It then applies conditional Flow Matching to the prediction-relative residual $R_{\hat B}=V-\hat B$ and composes $\hat V=\hat B+\hat R$. The same inference-available $\hat B$ is reused in residual supervision, residual conditioning, and final composition; we refer to this shared-reference construction as \emph{prediction consistency}.

The contributions of this study are threefold:
\begin{itemize}
    \item We develop role-aware conditioning for RMS velocity, post-stack migrated images, interpreted horizons, and sparse wells. All observed modalities can contribute to both roles, while the downstream interfaces explicitly distinguish background prediction from residual guidance.
    \item We formulate background-guided residual Flow Matching in which a separately supervised deterministic background defines the prediction-relative correction and conditions its generative transport, linking deterministic and generative reconstruction responsibilities.
    \item We enforce prediction consistency by reusing the same predicted background for residual supervision, residual conditioning, and final composition, and evaluate the resulting parameter-compact formulation under a common eight-subset OpenFWI-derived reconstruction protocol.
\end{itemize}

The resulting perspective is concise: \emph{predict the background; transport the prediction-relative correction}.
